\scanpage{22}\mspage{6 - 43 -}\hfill n° 392

e) Pour que $C^*$ soit pseudo-inversible, il faut et il suffit que pour tout idéal premier $\underline p$ de hauteur 1 de $A$, on ait $(\underline D_{C/A})_{\underline p}:(\underline D_{C/A})_{\underline p}=C_{\underline p}$.

4)\quad a) Soient $C$ un anneau noethérien intègre, $\unsure{\underline q|\deleted{\underline q}}$ un idéal primaire tel que $\underline q:\underline q=C$, $\underline p$ l'idéal premier correspondant ; on a alors $\underline q:\underline p\subset C$. Montrer que pour que l'idéal primaire $\underline q$ soit irréductible, il faut et il suffit que $(\underline q:\underline p)/\underline q$ soit un $C$-module simple. (Se ramener au cas où $C$ est un anneau artinien primaire avec $\underline q=0$, $\underline p$ étant donc l'unique idéal premier de $C$ ; montrer que $0:\underline p$ est le socle de $C$, en notant que c'est un $(A/\underline p)$-module, et pour que $0$ soit irréductible, il faut et il suffit que le socle de $C$ soit simple).

b) On suppose que les anneaux $A$, $B$, $C$ vérifient les hypothèses de l'exerc. 2. Montrer que, pour que l'idéal fractionnaire $C^*$ de $C$ soit pseudo-inversible, il faut et il suffit que pour tout idéal premier $\underline p$ de hauteur 1 dans $A$, les idéaux primaires de la décomposition primaire réduite, (unique), de $\underline p C_{\underline p}$ dans $C_{\underline p}$, soient irréductibles. (Se ramener au cas où $A$ est un anneau de valuation discrète d'uniformisante $\pi$ ; si $\underline m_i$, $(1\leq i\leq r)$, sont les idéaux maximaux de $C$, pour que $C^*$ soit inversible, il faut et il suffit que les $C$-modules $C^*/\underline m_i C^*$ soient simples, (cf. exerc. 1 b) ou, ce qui revient au même, que les $C$-modules $(\pi C)^*/\underline m_i(\pi C)^*$ soient simples ; utilisant la relation $(\underline a\underline b)^*=\underline a^*:\underline b$, en déduire que cette condition équivaut à dire que $(\pi C:\underline m_i)/\pi C$ est simple pour tout $i$ ; enfin, si $\underline q_i$ est la composante $\underline m_i$-primaire de $C$, observer que $\pi C:\underline m_i=\underline q_i:\underline m_i$ et utiliser a)).

c) Afin que, pour tout idéal premier $\underline p$ de hauteur 1 de $A$, les $C_{\underline p}$-modules $C^*_{\underline p}/C_{\underline p}$ et $C_{\underline p}/(\underline D_{C/A})_{\underline p}$ soient isomorphes, il faut et il suffit que l'idéal fractionnaire $C^*$ soit pseudo-inversible. (Se ramener au cas où $A$ est un anneau de valuation discrète ; observer que si $\underline m_i$, $(1\leq i\leq r)$, sont les idéaux maximaux de $C$, l'hypothèse que $C^*/C$ est isomorphe à $C/\underline D_{C/A}$ entraîne que $C^*/\underline m_i C^*$ est un $C$-module simple pour tout $i$).

d) Montrer que pour que l'on ait $C_{\underline p}=B_{\underline p}$ pour tout idéal premier $\underline p$ de hauteur 1 de $A$, il faut et il suffit que tout idéal primaire $\ne 0$ de $C$ soit pseudo-inversible. (Pour tout idéal premier $\underline p\ne0$ de $C$, considérer l'idéal $\underline p^2$).

5) Soit $K$ une extension algébrique finie de $Q$, $B$ la fermeture intégrale de $Z$ dans $K$. Afin que, pour tout sous-anneau $C$ de $B$ ayant même corps des fractions que $B$, $C^*$ soit un idéal fractionnaire inversible de $C$, il faut et il suffit que $[K:Q]\leq2$. (Si $[K:Q]=2$, montrer, en utilisant la prop. \quad, que $\underline D_{C/Z}$
% archive sheet number: 21
